% ch9_b (书页 389-400 / p404-p415)
%--A-TAIL--
%JOIN%并彼此正交归一，但具有负范数，

\begin{equation*}
(f^{*}_{\mathbf{k}_1},\, f^{*}_{\mathbf{k}_2}) = -\delta^{(n-1)}(\mathbf{k}_1-\mathbf{k}_2).
\tag{9.64}
\end{equation*}

模式 $f_{\mathbf{k}}$ 与 $f^{*}_{\mathbf{k}}$ 合起来构成一个完备集，借助它我们可以展开 Klein--Gordon 方程的任何解。

要对这一理论进行正则量子化，我们把经典变量（场及其共轭动量）提升为作用在希尔伯特空间上的算符，并在等时超曲面上施加正则对易关系：

\begin{gather*}
[\phi(t,\mathbf{x}),\, \phi(t,\mathbf{x}')] = 0 \\
[\pi(t,\mathbf{x}),\, \pi(t,\mathbf{x}')] = 0 \\
[\phi(t,\mathbf{x}),\, \pi(t,\mathbf{x}')] = i\,\delta^{(n-1)}(\mathbf{x}-\mathbf{x}').
\tag{9.65}
\end{gather*}

在场论中，我们需要显式地声明场与其动量在整个空间中与自身对易；对单个谐振子而言这是隐含的，因为那里只有一个坐标和一个动量，各自必然与自身对易。$\delta$ 函数意味着等时的算符在空间中处处对易，只在空间点重合处例外；这一特征源于因果性的要求（类空间隔的算符不能相互影响）。

正如 Klein--Gordon 方程的经典解可以用模式 (9.59) 展开，量子算符场 $\phi(t,\mathbf{x})$ 也可以。把场算符模式展开的系数记作 $\hat{a}^{\dagger}_{\mathbf{k}}$ 和 $\hat{a}_{\mathbf{k}}$，我们有

\begin{equation*}
\phi(t,\mathbf{x}) = \int d^{\,n-1}k\,\bigl[\hat{a}_{\mathbf{k}}f_{\mathbf{k}}(t,\mathbf{x}) + \hat{a}^{\dagger}_{\mathbf{k}}f^{*}_{\mathbf{k}}(t,\mathbf{x})\bigr].
\tag{9.66}
\end{equation*}

把这一展开代入式 (9.65)，我们发现算符 $\hat{a}^{\dagger}_{\mathbf{k}}$ 和 $\hat{a}_{\mathbf{k}}$ 满足对易关系

\begin{gather*}
[\hat{a}_{\mathbf{k}},\, \hat{a}_{\mathbf{k}'}] = 0 \\
[\hat{a}^{\dagger}_{\mathbf{k}},\, \hat{a}^{\dagger}_{\mathbf{k}'}] = 0 \\
[\hat{a}_{\mathbf{k}},\, \hat{a}^{\dagger}_{\mathbf{k}'}] = \delta^{(n-1)}(\mathbf{k}-\mathbf{k}').
\tag{9.67}
\end{gather*}

于是这些算符满足产生和湮灭算符特有的对易关系，与简单谐振子的式 (9.23) 类似。当然，区别在于这样的算符有无穷多个，以 $\mathbf{k}$ 标记。由此可以看出把模式分为正频与负频的意义所在：正频模式是湮灭算符的系数，而负频模式是产生算符的系数。正、负频模式的概念以后会推广到静态时空，但不会推广到任意时空。

在谐振子的情形中，我们曾用产生和湮灭算符为希尔伯特空间定义了一组基，其中的基矢是粒子数算符（number operator）的本征态。同样的步骤也适用于自由标量场，只是现在我们必须对每个空间波矢 $\mathbf{k}$ 分别记录激发的数目。将会有一个唯一的真空态 $|0\rangle$，其特征是它能被每一个 $\hat{a}_{\mathbf{k}}$ 湮灭，

\begin{equation*}
\hat{a}_{\mathbf{k}}|0\rangle = 0 \qquad \text{for all } \mathbf{k}.
\tag{9.68}
\end{equation*}

具有 $n_{\mathbf{k}}$ 个相同动量 $\mathbf{k}$ 粒子的态，由 $\hat{a}^{\dagger}_{\mathbf{k}}$ 的反复作用产生，

\begin{equation*}
|n_{\mathbf{k}}\rangle = \frac{1}{\sqrt{n_{\mathbf{k}}!}}\left(\hat{a}^{\dagger}_{\mathbf{k}}\right)^{n_{\mathbf{k}}}|0\rangle,
\tag{9.69}
\end{equation*}

而具有各种动量 $\mathbf{k}_i$ 的 $n_i$ 个激发的态则为

\begin{equation*}
|n_1, n_2, \ldots, n_j\rangle = \frac{1}{\sqrt{n_1!n_2!\cdots n_j!}}\left(\hat{a}^{\dagger}_{\mathbf{k}_1}\right)^{n_1}\left(\hat{a}^{\dagger}_{\mathbf{k}_2}\right)^{n_2}\cdots\left(\hat{a}^{\dagger}_{\mathbf{k}_j}\right)^{n_j}|0\rangle.
\tag{9.70}
\end{equation*}

作用在这样的态上，产生和湮灭算符会如预期那样改变激发的数目：

\begin{align*}
\hat{a}_{\mathbf{k}_i}|n_1, n_2, \ldots, n_i, \ldots, n_j\rangle &= \sqrt{n_i}\,|n_1, n_2, \ldots, n_i-1, \ldots, n_j\rangle \\
\hat{a}^{\dagger}_{\mathbf{k}_i}|n_1, n_2, \ldots, n_i, \ldots, n_j\rangle &= \sqrt{n_i+1}\,|n_1, n_2, \ldots, n_i+1, \ldots, n_j\rangle.
\tag{9.71}
\end{align*}

我们可以对每个波矢定义一个粒子数算符，

\begin{equation*}
\hat{n}_{\mathbf{k}} = \hat{a}^{\dagger}_{\mathbf{k}}\hat{a}_{\mathbf{k}},
\tag{9.72}
\end{equation*}

它满足

\begin{equation*}
\hat{n}_{\mathbf{k}_i}|n_1, n_2, \ldots, n_i, \ldots, n_j\rangle = n_i|n_1, n_2, \ldots, n_i, \ldots, n_j\rangle.
\tag{9.73}
\end{equation*}

这些粒子数算符的本征态构成整个希尔伯特空间的一组基，称为 \textbf{Fock 基}（Fock basis）；由这组基构造的空间常被称作“Fock 空间”，但它当然就是原来的希尔伯特空间。

我们可能想考察的一件事，是我们的 Fock 基在洛伦兹变换下的行为。我们显然一直在利用 Minkowski 空间的对称性，例如用平面波作为 Klein--Gordon 方程解的一组基。这些模式的关键特征在于我们能够区分正频与负频，从而可以把 $\phi$ 的模式展开中的系数解释为湮灭和产生算符。现在考虑速度为 $\mathbf{v} = d\mathbf{x}/dt$ 的推动（boost），它给出新坐标 $x^{\mu'}$：

\begin{equation*}
t' = \gamma t - \gamma\,\mathbf{v}\cdot\mathbf{x}, \qquad \mathbf{x}' = \gamma\,\mathbf{x} - \gamma\,\mathbf{v}t,
\tag{9.74}
\end{equation*}

其中 $\gamma = 1/\sqrt{1-v^2}$，逆变换则为

\begin{equation*}
t = \gamma t' + \gamma\,\mathbf{v}\cdot\mathbf{x}', \qquad \mathbf{x} = \gamma\,\mathbf{x}' + \gamma\,\mathbf{v}t'.
\tag{9.75}
\end{equation*}

我们的模式函数在推动后的参考系中的时间导数为

\begin{align*}
\partial_{t'}f_{\mathbf{k}} &= \frac{\partial x^{\mu}}{\partial t'}\partial_{\mu}f_{\mathbf{k}} \\
&= \gamma(-i\omega)f_{\mathbf{k}} + \gamma\,\mathbf{v}\cdot(i\mathbf{k})f_{\mathbf{k}} \\
&= -i\omega' f_{\mathbf{k}}
\tag{9.76}
\end{align*}

其中

\begin{equation*}
\omega' = \gamma\omega - \gamma\,\mathbf{v}\cdot\mathbf{k}
\tag{9.77}
\end{equation*}

不过是在推动后的参考系中的频率。于是很明显，一个描述具有某些动量的粒子集合的态，被推动成描述同样一些粒子、但动量经过推动的态。因此，两个参考系中的总粒子数算符相一致，特别是真空态也相一致。在这个意义上，我们最初对惯性系的选择是无关紧要的。在下一节中我们将看到，我们之所以能找到正频和负频解，可以追溯到 Minkowski 时空中类时 Killing 矢量 $\partial_t$ 的存在；而 Fock 空间在基变换下的不变性，则可以追溯到所有这些类时 Killing 矢量都由洛伦兹变换相联系这一事实。因此，即使一个模式的频率依赖于惯性系的选择，正负频的分解却是不变的。

我们希望把哈密顿量

\begin{equation*}
H = \int d^{\,n-1}x\left[\frac{1}{2}\dot{\phi}^2 + \frac{1}{2}(\nabla\phi)^2 + \frac{1}{2}m^2\phi^2\right]
\tag{9.78}
\end{equation*}

用产生和湮灭算符表示出来，就像我们对谐振子做过的那样。我们可以逐项分析这个表达式，为简单起见从 $\phi^2$ 项开始：

\begin{align*}
&\frac{1}{2}m^2\int d^{\,n-1}x\,\phi^2 \\
&\quad = \frac{1}{2}m^2\int d^{\,n-1}x\, d^{\,n-1}k\, d^{\,n-1}k'\Bigl(\hat{a}_{\mathbf{k}}f_{\mathbf{k}} + \hat{a}^{\dagger}_{\mathbf{k}}f^{*}_{\mathbf{k}}\Bigr)\Bigl(\hat{a}_{\mathbf{k}'}f_{\mathbf{k}'} + \hat{a}^{\dagger}_{\mathbf{k}'}f^{*}_{\mathbf{k}'}\Bigr) \\
&\quad = \frac{1}{2}m^2\int d^{\,n-1}x\, d^{\,n-1}k\, d^{\,n-1}k'\Bigl(\hat{a}_{\mathbf{k}}\hat{a}_{\mathbf{k}'}f_{\mathbf{k}}f_{\mathbf{k}'} + \hat{a}^{\dagger}_{\mathbf{k}}\hat{a}_{\mathbf{k}'}f^{*}_{\mathbf{k}}f_{\mathbf{k}'} \\
&\qquad + \hat{a}_{\mathbf{k}}\hat{a}^{\dagger}_{\mathbf{k}'}f_{\mathbf{k}}f^{*}_{\mathbf{k}'} + \hat{a}^{\dagger}_{\mathbf{k}}\hat{a}^{\dagger}_{\mathbf{k}'}f^{*}_{\mathbf{k}}f^{*}_{\mathbf{k}'}\Bigr).
\tag{9.79}
\end{align*}

细看括号中的第一项，并暂时忽略对 $\mathbf{k}$ 的积分，代入模式函数 (9.59) 的显式形式，可得

\begin{align*}
\int d^{\,n-1}x\, d^{\,n-1}k'\,\hat{a}_{\mathbf{k}}\hat{a}_{\mathbf{k}'}f_{\mathbf{k}}f_{\mathbf{k}'}
&= \int d^{\,n-1}x\, d^{\,n-1}k'\,\hat{a}_{\mathbf{k}}\hat{a}_{\mathbf{k}'}\frac{e^{-i(\omega+\omega')t}e^{i(\mathbf{k}+\mathbf{k}')\cdot\mathbf{x}}}{2(2\pi)^{n-1}\sqrt{\omega\omega'}} \\
&= \int d^{\,n-1}k'\,\hat{a}_{\mathbf{k}}\hat{a}_{\mathbf{k}'}\frac{e^{-i(\omega+\omega')t}}{2\sqrt{\omega\omega'}}\,\delta^{(n-1)}(\mathbf{k}+\mathbf{k}') \\
&= \hat{a}_{\mathbf{k}}\hat{a}_{-\mathbf{k}}\frac{e^{-2i\omega t}}{2\omega},
\tag{9.80}
\end{align*}

这里我们再次用到了式 (9.58)。类似地计算式 (9.79) 中的其他项，我们发现对哈密顿量的势能贡献于是变为

\begin{align*}
\frac{1}{2}m^2\int d^{\,n-1}x\,\phi^2 = \frac{1}{2}m^2\int d^{\,n-1}k\left(\frac{1}{2\omega}\right)\Bigl[&\hat{a}_{\mathbf{k}}\hat{a}_{-\mathbf{k}}e^{-2i\omega t} \\
&+ \hat{a}^{\dagger}_{\mathbf{k}}\hat{a}_{\mathbf{k}} + \hat{a}_{\mathbf{k}}\hat{a}^{\dagger}_{\mathbf{k}} + \hat{a}^{\dagger}_{\mathbf{k}}\hat{a}^{\dagger}_{-\mathbf{k}}e^{2i\omega t}\Bigr].
\tag{9.81}
\end{align*}

对于动能和梯度能这两部分，求导会分别拉下 $\omega$ 和 $\mathbf{k}$ 的因子；我们得到

\begin{equation*}
\frac{1}{2}\int d^{\,n-1}x\,\dot{\phi}^2 = \frac{1}{2}\int d^{\,n-1}k\left(\frac{\omega}{2}\right)\Bigl[-\hat{a}_{\mathbf{k}}\hat{a}_{-\mathbf{k}}e^{-2i\omega t} + \hat{a}^{\dagger}_{\mathbf{k}}\hat{a}_{\mathbf{k}} + \hat{a}_{\mathbf{k}}\hat{a}^{\dagger}_{\mathbf{k}} - \hat{a}^{\dagger}_{\mathbf{k}}\hat{a}^{\dagger}_{-\mathbf{k}}e^{2i\omega t}\Bigr]
\tag{9.82}
\end{equation*}

以及

\begin{align*}
\frac{1}{2}\int d^{\,n-1}x\,(\nabla\phi)^2 = \frac{1}{2}\int d^{\,n-1}k\left(\frac{\mathbf{k}^2}{2\omega}\right)\Bigl[&\hat{a}_{\mathbf{k}}\hat{a}_{-\mathbf{k}}e^{-2i\omega t} \\
&+ \hat{a}^{\dagger}_{\mathbf{k}}\hat{a}_{\mathbf{k}} + \hat{a}_{\mathbf{k}}\hat{a}^{\dagger}_{\mathbf{k}} + \hat{a}^{\dagger}_{\mathbf{k}}\hat{a}^{\dagger}_{-\mathbf{k}}e^{2i\omega t}\Bigr].
\tag{9.83}
\end{align*}

利用 $\omega^2 = \mathbf{k}^2 + m^2$，我们可以把所有结果合在一起，把标量场理论的哈密顿量写成

\begin{align*}
H &= \frac{1}{2}\int d^{\,n-1}k\,\bigl[\hat{a}^{\dagger}_{\mathbf{k}}\hat{a}_{\mathbf{k}} + \hat{a}_{\mathbf{k}}\hat{a}^{\dagger}_{\mathbf{k}}\bigr]\omega \\
&= \int d^{\,n-1}k\,\Bigl[\hat{n}_{\mathbf{k}} + \frac{1}{2}\delta^{(n-1)}(0)\Bigr]\omega,
\tag{9.84}
\end{align*}

其中最后一步用到了对易关系 (9.67) 以及粒子数算符 $\hat{n}_{\mathbf{k}} = \hat{a}^{\dagger}_{\mathbf{k}}\hat{a}_{\mathbf{k}}$。按类似的逻辑，我们可以构造一个对应于总动量空间分量的算符，算出来是

\begin{equation*}
P^i = \int d^{\,n-1}k\,\hat{n}_{\mathbf{k}}k^i.
\tag{9.85}
\end{equation*}

正如我们所预期的，能量本征态将是具有固定激发数目的那些态，每个激发携带能量 $\omega$。Fock 基中的激发被解释为粒子。这就是粒子在量子场论中产生的方式：能量本征态是具有确定动量的粒子集合。当然，我们的模式是弥漫整个空间的平面波，而不是我们想到粒子时脑海中浮现的那种气泡室里的局域径迹。更糟的是，在弯曲时空中，波动方程将不再有我们能解释为粒子的、具有确定频率的平面波解。解决这两个问题的办法是操作性地思考，即从实验装置会观测到什么出发。最好的策略是先定义一个合理的粒子探测器（particle detector）概念，它在平直时空中还原为我们的直观图像，然后把“粒子”定义为“粒子探测器探测到的东西”。对一个定义恰当的粒子探测器，可以证明我们的平面波模式会像我们希望的那样“留下径迹”；在这样的探测器阵列中，如果一个平面波触发了一个探测器，那么它沿着从第一个探测器出发、由波矢给出的方向的路径，触发其他探测器的概率很高。（我们应该指出，如果你到 Fermilab 或 CERN 这样的地方参观一座真正的粒子加速器，会发现那些探测器与研究弯曲时空量子场论的理论家们所发明的探测器几乎毫无相似之处；不过究其本质，二者存在根本的相似性。）关于粒子探测器的讨论可参见 Birrell 和 Davies (1982)。

你或许会为哈密顿量 (9.84) 中的因子 $\delta^{(n-1)}(0)$ 感到担心——而且你理应如此。它意味着即使在真空态 $|0\rangle$ 下测量，哈密顿量也是无穷大。这一项是谐振子零点能 (9.20) 在场论中的对应物。在第 4 章讨论宇宙学常数时我们提到过，量子涨落会引起经典真空能形式上无穷大的位移；当把引力包括进来时，标量场哈密顿量中的这个无穷大贡献将表现为一个发散的宇宙学常数。它是无穷大量 $\delta^{(n-1)}(0)$ 在无穷大的 $\mathbf{k}$ 范围上的积分这一事实，可以转述为：总能量是对无限大空间中无穷大能量密度的积分。但真正的问题在于高频模式贡献的能量密度，而不是无限大的体积；如果我们把计算放进体积为 $L^{n-1}$ 的盒子里来做正则化，就会发现

\begin{equation*}
\frac{1}{2}\int d^{\,n-1}k\,\delta^{(n-1)}(0)\,\omega \to \frac{1}{2}\left(\frac{L}{2\pi}\right)^{n-1}\sum_{\mathbf{k}}\omega,
\tag{9.86}
\end{equation*}

即使 $L$ 有限它也发散，因为 $\mathbf{k}$（从而 $\omega$）可以任意大。在某个高动量 $k_{\mathrm{max}}$ 处放一个截断，就会回到式 (4.104)。

在简单谐振子的情形中我们指出过，当初要是选一个具有负极小值的经典势，零点能本可以避免；量子力学的贡献不一定代表真正的答案，只代表能量相对其经典值的位移。场论中也是如此；我们可以自由地定义最初的经典标量场理论，使得量子力学真空能为零。然而，我们不能逐模式地减去有限的能量，因为我们的自由仅仅是往势里、从而往哈密顿量密度 (9.52) 中加一个单独的常数。要得到真空态的有限哈密顿量，这个常数必须是无穷大。减去一个无穷大常数并没有什么不对；这是量子场论中一种历史悠久的技术，称为“重整化”（renormalization）。重整化有时看上去吓人，或者多少有些不正当，但实际上它完全合理；无穷大只出现在量子理论与它们经典对应物之间的关系之中，而不出现在任何可观测量之中。既然大自然多半既不知道也不在乎我们对经典力学的偏爱，重整化就不应该有什么特别令人不安的地方。

当然，一旦我们重整化得到有限的真空能，这个能量就可以取我们喜欢的任何值；它完全任意。这一点在弯曲时空量子场论中同样成立；我们也许无法把场分解成具有确定频率的模式，因而无法给每个模式指派一个真空能贡献，但仔细的分析允许我们把真空能重整化成你喜欢的任何数。同样，并没有发生什么深刻的事情；真空能在我们的经典模型中本来就是完全任意的，我们只是为了方便才把它选成零。宇宙学常数问题并不是因为量子力学贡献了大量真空能而产生的，因为这个贡献可以直接重整化掉；问题在于，所得到的任意数没有理由接近零。如前所述，从有效场论的角度看，问题还要更尖锐一些，因为对真空能的尺度存在一个逻辑上的预期，那就是未知的量子引力效应应当开始起作用的 Planck 标度。不过，贯穿本章我们只关心量子场在固定时空背景中的传播，而不把量子能动张量用作 Einstein 方程的源；因此我们可以选择无视宇宙学常数问题。

\section{弯曲时空中的量子场论}

在第 4 章中我们讨论过，把物理理论从平直时空推广到弯曲时空是多么容易——我们只需把理论写成坐标不变的形式，然后断言当时空弯曲时它们仍然成立。这一步骤对量子场论依然有效，尽管我们将不得不放弃一些在平直时空中显得不可或缺的概念。

我们从弯曲时空中标量场的拉格朗日密度出发，

\begin{equation*}
\mathcal{L} = \sqrt{-g}\left(-\frac{1}{2}g^{\mu\nu}\nabla_{\mu}\phi\nabla_{\nu}\phi - \frac{1}{2}m^2\phi^2 - \xi R\phi^2\right).
\tag{9.87}
\end{equation*}

除了可以预料到的度规 $g_{\mu\nu}$ 及其行列式的出现之外，我们还包含了与标量曲率 $R$ 的直接耦合，由一个常数 $\xi$ 参数化。由于 $\xi$ 无量纲，没有理由预期它很小；实际上，它自然应该是 1 的量级。文献中对 $\xi$ 的值有两种受青睐的选择：\textbf{最小耦合}（minimal coupling）干脆关掉与 $R$ 的直接相互作用，

\begin{equation*}
\xi = 0,
\tag{9.88}
\end{equation*}

而\textbf{共形耦合}（conformal coupling）则取

\begin{equation*}
\xi = \frac{(n-2)}{4(n-1)},
\tag{9.89}
\end{equation*}

在四维时它就是 $\xi = \frac{1}{6}$。利用附录 G 中的公式容易验证，当 $\xi$ 取这个值且 $m=0$ 时，标量场理论在共形变换 $g_{\mu\nu}\to\omega^2(x)g_{\mu\nu}$ 下不变。实际上，在真实世界中并没有充分的理由在最小耦合和共形耦合之间择取其一；最小耦合不增强任何对称性，而共形不变性显然不是大多数物理理论所具有的对称性。（由于共形变换是局部的标度改变，以质量这样的有量纲参数为特征的理论一般不会共形不变。）即使经典理论是共形不变的，量子化也可能破坏这一对称性，例如在与无质量夸克耦合的量子色动力学（QCD）理论中就发生了这种情况。一般而言，在四维中很难找到严格共形不变的相互作用理论，尽管已知某些具有高度超对称性的模型是共形不变的。

我们可以像先前那样继续量子化这个理论。共轭动量是

\begin{equation*}
\pi = \frac{\partial\mathcal{L}}{\partial(\nabla_0\phi)},
\tag{9.90}
\end{equation*}

对于拉格朗日量 (9.87)，它为

\begin{equation*}
\pi = \sqrt{-g}\,\nabla_0\phi.
\tag{9.91}
\end{equation*}

我们可以施加正则对易关系

\begin{gather*}
[\phi(t,\mathbf{x}),\, \phi(t,\mathbf{x}')] = 0 \\
[\pi(t,\mathbf{x}),\, \pi(t,\mathbf{x}')] = 0 \\
[\phi(t,\mathbf{x}),\, \pi(t,\mathbf{x}')] = \frac{i}{\sqrt{-g}}\,\delta^{(n-1)}(\mathbf{x}-\mathbf{x}').
\tag{9.92}
\end{gather*}

标量场的运动方程是

\begin{equation*}
\Box\phi - m^2\phi - \xi R\phi = 0.
\tag{9.93}
\end{equation*}

对具有诱导度规 $\gamma_{ij}$ 和单位法矢量 $n^{\mu}$ 的类空超曲面 $\Sigma$，该方程解上的内积为

\begin{equation*}
(\phi_1, \phi_2) = -i\int_{\Sigma}\left(\phi_1\nabla_{\mu}\phi^{*}_2 - \phi^{*}_2\nabla_{\mu}\phi_1\right)n^{\mu}\sqrt{\gamma}\, d^{\,n-1}x,
\tag{9.94}
\end{equation*}

它与 $\Sigma$ 的选取无关。

到目前为止一切顺利。为了继续我们在平直空间中所采取的步骤，现在应当引进一组正频和负频模式，它们构成式 (9.93) 解的完备基，把场算符 $\phi$ 用这些模式展开，并把算符系数解释为产生和湮灭算符。正是在这一步，我们的程序失灵了。由于一般不存在任何类时 Killing 矢量，我们通常无法找到能分离成时间依赖和空间依赖因子的波动方程的解，相应地也就不能把模式分类为正频或负频。我们可以找到一组基模式，但问题在于这样的模式组一般有很多，没有任何办法偏爱其中一组而舍弃其他，而真空和粒子数算符的概念将敏感地依赖于我们选择哪一组。

看看我们能做些什么。我们总能找到式 (9.93) 的一组正交归一的解 $f_i(x^{\mu})$，

\begin{equation*}
(f_i, f_j) = \delta_{ij},
\tag{9.95}
\end{equation*}

以及具有负范数的相应共轭模式，

\begin{equation*}
(f^{*}_i, f^{*}_j) = -\delta_{ij}.
\tag{9.96}
\end{equation*}

指标 $i$ 可以连续或分立；目前我们采用适合分立情形的记号。这些模式可以选成完备集，于是我们可以把场展开为

\begin{equation*}
\phi = \sum_i\left(\hat{a}_i f_i + \hat{a}^{\dagger}_i f^{*}_i\right).
\tag{9.97}
\end{equation*}

系数 $\hat{a}_i$ 和 $\hat{a}^{\dagger}_i$ 满足对易关系

\begin{gather*}
[\hat{a}_i,\, \hat{a}_j] = 0 \\
[\hat{a}^{\dagger}_i,\, \hat{a}^{\dagger}_j] = 0 \\
[\hat{a}_i,\, \hat{a}^{\dagger}_j] = \delta_{ij}.
\tag{9.98}
\end{gather*}

将存在一个真空态 $|0_f\rangle$，它被所有湮灭算符湮灭，

\begin{equation*}
\hat{a}_i|0_f\rangle = 0 \qquad \text{for all } i.
\tag{9.99}
\end{equation*}

从这个真空态出发，我们可以为希尔伯特空间定义整套 Fock 基。与之前一样，具有 $n_i$ 个激发的态由 $\hat{a}^{\dagger}_i$ 的反复作用产生，

\begin{equation*}
|n_i\rangle = \frac{1}{\sqrt{n_i!}}\left(\hat{a}^{\dagger}_i\right)^{n_i}|0_f\rangle,
\tag{9.100}
\end{equation*}

对具有不同种类激发的态也同样如此。我们甚至可以对每个模式定义一个粒子数算符，

\begin{equation*}
\hat{n}_{fi} = \hat{a}^{\dagger}_i\hat{a}_i.
\tag{9.101}
\end{equation*}

真空态和粒子数算符下标中的 $f$ 提醒我们，它们是相对于模式组 $f_i$ 定义的。

这套架构看上去与我们在平直空间中的非常相似；为什么我们不能宣布由 $\hat{a}^{\dagger}_i$ 产生的激发就是粒子，然后就此收工呢？我们本可以这样做，但必须面对这样一个事实：还存在其他可能的选择；基模式 $f_i(x^{\mu})$ 高度不唯一。考虑另外一组模式 $g_i(x^{\mu})$，它具有我们原来的模式所具有的全部性质，包括（连同共轭模式 $g^{*}_i$）构成一个完备基，相对于它可以展开我们的场算符，

\begin{equation*}
\phi = \sum_i\left(\hat{b}_i g_i + \hat{b}^{\dagger}_i g^{*}_i\right).
\tag{9.102}
\end{equation*}

湮灭和产生算符 $\hat{b}_i$ 与 $\hat{b}^{\dagger}_i$ 满足对易关系

\begin{gather*}
[\hat{b}_i,\, \hat{b}_j] = 0 \\
[\hat{b}^{\dagger}_i,\, \hat{b}^{\dagger}_j] = 0 \\
[\hat{b}_i,\, \hat{b}^{\dagger}_j] = \delta_{ij},
\tag{9.103}
\end{gather*}

并且将存在一个被所有湮灭算符湮灭的真空态 $|0_g\rangle$，

\begin{equation*}
\hat{b}_i|0_g\rangle = 0 \qquad \text{for all } i.
\tag{9.104}
\end{equation*}

我们可以通过在这个真空上反复作用产生算符来构造 Fock 基，并定义粒子数算符

\begin{equation*}
\hat{n}_{gi} = \hat{b}^{\dagger}_i\hat{b}_i.
\tag{9.105}
\end{equation*}

在从平直时空到弯曲时空的转变中，我们失去的是偏爱某一组模式甚于其他任何一组的任何理由。在平直时空中，我们能够通过要求模式相对于时间坐标为正频〔如式 (9.61) 所定义的〕而挑出一组自然的模式。时间坐标并不唯一，因为我们可以自由地做洛伦兹变换；但我们看到过，真空态和总粒子数算符在这样的变换下不变。因此，每个惯性系观者对什么是真空态、周围有多少粒子，都会有一致的看法。

在我们现在考虑的更一般情形中，如果一个观者相对于模式组 $f_i$ 定义粒子，而另一个观者使用模式组 $g_i$，那么他们对观测到多少个粒子（甚至究竟有没有观测到粒子）通常会有分歧。为了看清这一点，方便的做法是把每组模式用另一组展开，

\begin{align*}
g_i &= \sum_j\left(\alpha_{ij}f_j + \beta_{ij}f^{*}_j\right) \\
f_i &= \sum_j\left(\alpha^{*}_{ji}g_j - \beta_{ji}g^{*}_j\right).
\tag{9.106}
\end{align*}

从一组基模式到另一组的变换称为 \textbf{Bogoliubov 变换}（Bogoliubov transformation），实现该变换的矩阵 $\alpha_{ij}$ 和 $\beta_{ij}$ 称为 Bogoliubov 系数。利用模式函数的正交归一性，它们可以表示为

\begin{gather*}
\alpha_{ij} = (g_i, f_j) \\
\beta_{ij} = -(g_i, f^{*}_j).
\tag{9.107}
\end{gather*}

它们满足自己的归一化条件，

\begin{align*}
\sum_j\left(\alpha_{ik}\alpha^{*}_{jk} - \beta_{ik}\beta^{*}_{jk}\right) &= \delta_{ij} \\
\sum_j\left(\alpha_{ik}\beta_{jk} - \beta_{ik}\alpha_{jk}\right) &= 0.
\tag{9.108}
\end{align*}

Bogoliubov 系数除了描述模式之间的变换，还可以用来在算符之间变换，

\begin{align*}
\hat{a}_i &= \sum_j\left(\alpha_{ji}\hat{b}_j + \beta^{*}_{ji}\hat{b}^{\dagger}_j\right) \\
\hat{b}_i &= \sum_j\left(\alpha^{*}_{ij}\hat{a}_j - \beta^{*}_{ij}\hat{a}^{\dagger}_j\right).
\tag{9.109}
\end{align*}

现在设想系统处于 $f$ 真空 $|0_f\rangle$ 中，其中不会观测到任何 $f$ 粒子；我们想知道使用 $g$ 模式的观者会观测到多少粒子。因此我们来计算 $g$ 粒子数算符在 $f$ 真空中的期望值：

\begin{align*}
\langle 0_f|\hat{n}_{gi}|0_f\rangle &= \langle 0_f|\hat{b}^{\dagger}_i\hat{b}_i|0_f\rangle \\
&= \left\langle 0_f\left|\sum_{jk}\left(\alpha_{ij}\hat{a}^{\dagger}_j - \beta_{ij}\hat{a}_j\right)\left(\alpha^{*}_{ik}\hat{a}_k - \beta^{*}_{ik}\hat{a}^{\dagger}_k\right)\right|0_f\right\rangle \\
&= \sum_{jk}\left(-\beta_{ij}\right)\left(-\beta^{*}_{ik}\right)\langle 0_f|\hat{a}_j\hat{a}^{\dagger}_k|0_f\rangle \\
&= \sum_{jk}\beta_{ij}\beta^{*}_{ik}\langle 0_f|\left(\hat{a}^{\dagger}_k\hat{a}_j + \delta_{jk}\right)|0_f\rangle \\
&= \sum_{jk}\beta_{ij}\beta^{*}_{ik}\delta_{jk}\langle 0_f|0_f\rangle \\
&= \sum_j\beta_{ij}\beta^{*}_{ij}.
\tag{9.110}
\end{align*}

于是，$f$ 真空中 $g$ 粒子的数目可以用 Bogoliubov 系数表示为

\begin{equation*}
\langle 0_f|\hat{n}_{gi}|0_f\rangle = \sum_j|\beta_{ij}|^2.
\tag{9.111}
\end{equation*}

没有理由认为它为零；从一个视角看是空无一物的真空，从另一个视角看却翻腾着粒子。只要任何一个 $\beta_{ij}$ 非零，两个真空态就不会重合。为什么会这样，看看式 (9.109) 就能明白：$\beta_{ij}$ 描述的是一个基中的产生算符混入另一个基中湮灭算符的成分。

这些关于模式和粒子数算符的讨论可能显得不必要地抽象；诚然，如果一个真正的粒子探测器在可能是弯曲的时空中沿某条轨迹运动，它要么探测到粒子，要么探测不到，并不知道我们的场论用的是哪一组基模式。我们怎么知道这样的探测器实际使用的是哪一种“粒子”定义？答案在于：探测器测量的是沿其轨迹的固有时 $\tau$，并且将相对于该固有时来定义正频和负频。因此，如果能找到一组满足

\begin{equation*}
\frac{D}{d\tau}f_i = -i\omega f_i,
\tag{9.112}
\end{equation*}

的模式 $f_i$，我们就能用这些模式计算探测器将看到多少粒子。当然，一般不可能在整个时空中都找到这样的模式。有可能做到的一种情形是\emph{静态}（static）时空，此时我们有一个超曲面正交的类时 Killing 矢量 $K^{\mu}$。在那种情况下，我们可以选取这样的坐标：度规分量与时间坐标 $t$ 无关，并且不存在时间-空间交叉项：

\begin{equation*}
\partial_0 g_{\mu\nu} = 0, \qquad g_{0i} = 0.
\tag{9.113}
\end{equation*}

（指标 $i$、$j$ 现在指空间分量，不是模式标记。）对这样的度规，d'Alembert 算符作用在某个模式函数 $f(t,\mathbf{x})$ 上的结果是

\begin{equation*}
\Box f = \left[g^{00}\partial_0^2 + \frac{1}{2}g^{00}g^{ij}(\partial_i g_{00})\partial_j + g^{ij}\partial_i\partial_j - g^{ij}\Gamma^{k}_{ij}\partial_k\right]f.
\tag{9.114}
\end{equation*}

于是运动方程 (9.93) 可以写成

\begin{equation*}
\partial_0^2 f = -\left(g^{00}\right)^{-1}\left[g^{ij}\partial_i\partial_j + \frac{1}{2}g^{00}g^{ij}(\partial_i g_{00})\partial_j - g^{ij}\Gamma^{k}_{ij}\partial_k - \left(m^2 + \xi R\right)\right]f.
\tag{9.115}
\end{equation*}

左边的算符是纯时间导数，而右边的算符只涉及空间导数和仅依赖空间的函数。因此我们可以找到分离变量的解

\begin{equation*}
f_{\omega}(t,\mathbf{x}) = e^{-i\omega t}\bar{f}_{\omega}(\mathbf{x}),
\tag{9.116}
\end{equation*}

它们可以称为正频的，

\begin{equation*}
\partial_t f_{\omega}(t,\mathbf{x}) = -i\omega f_{\omega}(t,\mathbf{x}), \qquad \omega > 0.
\tag{9.117}
\end{equation*}

这个关系可以改写成坐标不变的形式

\begin{equation*}
\mathcal{L}_K f_{\omega} = K^{\mu}\partial_{\mu}f_{\omega} = -i\omega f_{\omega}, \qquad \omega > 0,
\tag{9.118}
\end{equation*}

其中 $\mathcal{L}_K f_{\omega}$ 表示 $f_{\omega}$ 沿 $K$ 的李导数。同时也存在负频的共轭模式，

\begin{equation*}
\mathcal{L}_K f^{*}_{\omega} = K^{\mu}\partial_{\mu}f^{*}_{\omega} = i\omega f^{*}_{\omega}, \qquad \omega > 0.
\tag{9.119}
\end{equation*}

模式 $(f_{\omega}, f^{*}_{\omega})$ 合起来将构成静态背景中波动方程解的一组基。除非这些模式与我们的探测器相关，否则它们的存在帮不了我们；如果探测器的轨迹沿着 Killing 场的轨道行进（四速度 $U^{\mu} = dx^{\mu}/d\tau$ 与 $K^{\mu}$ 成正比），固有时就将与 Killing 时间 $t$ 成正比，而相对于这个 Killing 矢量为正频的模式便可以作为描述 Fock 空间的自然基。在下一节讨论安鲁效应时，我们将看到这一现象的实际运作。

在上一节中我们提到过，量子场论中需要对真空能进行重整化。这一需求在弯曲时空中依然存在，但由于不存在偏好的模式基，构造恰当的重整化程序更加困难。不过，至少在某些情形下，人们已经发展出代数方法来严格地定义重整化的能动张量；我们不会深入细节地探讨这个专题，但至少应当介绍一下其背后的基本思想。基本想法是：即使存在曲率，在足够小的尺度上时空看上去应当是 Minkowski 式的。因为我们在平直时空中发现的真空能发散源于短波长模式，所以我们应当能把弯曲时空中场在非常小尺度上的行为与平直时空中的行为相匹配，并减去出现的任何发散。特别地，我们考虑量子场 $\phi$ 在某个态 $|\psi\rangle$ 中的两点函数，
